% ============================================================
% introduzione_relazione.tex
% ============================================================
\section*{Introduzione}

Lo shopping online su marketplace come eBay presenta un problema 
fondamentale: la quantità di prodotti disponibili rende la ricerca 
tradizionale per parole chiave insufficiente. L'utente si trova a 
dover valutare manualmente l'affidabilità dei venditori, confrontare 
prezzi su inserzioni eterogenee e districarsi tra prodotti principali, 
accessori e ricambi spesso mescolati tra loro.\\

\noindent
Il progetto \textit{MCP E-commerce AI Shopping Assistant} nasce per 
automatizzare questo processo. Il sistema espone una chat 
conversazionale in cui l'utente formula richieste in linguaggio 
naturale — opzionalmente accompagnate da immagini — e riceve 
raccomandazioni motivate in tempo reale: l'agente si occupa 
autonomamente di interpretare la richiesta, pianificare le operazioni 
necessarie, interrogare il marketplace e sintetizzare i risultati.\\

\noindent
L'architettura è costruita attorno al \textit{Model Context Protocol} 
(MCP), standard aperto di Anthropic che disaccoppia la logica 
dell'agente dai servizi esterni. Il backend è sviluppato in Python con 
FastAPI; il frontend in React/TypeScript. La persistenza è affidata a 
PostgreSQL, Redis e Qdrant (database vettoriale). Le funzionalità di 
linguaggio e visione artificiale si appoggiano a Ollama Cloud, mentre 
i dati di prodotto provengono dalle API ufficiali di eBay. 
L'automazione browser è gestita tramite Playwright.\\

\noindent
I capitoli che seguono descrivono nel dettaglio i componenti 
principali: il protocollo MCP e l'agente ReAct (Capitolo 1), la pipeline 
di elaborazione della query (Capitolo 2), il sistema di ricerca e ranking 
(Capitolo 3), l'analisi dei venditori (Capitolo 4), gli strumenti specializzati 
(Capitolo 5), la personalizzazione e la memoria (Capitolo 6) e l'infrastruttura 
(Capitolo 7).